\documentclass[a4paper, 11pt]{article}
\usepackage[utf8]{inputenc}
\usepackage[T1]{fontenc}

\usepackage{amsmath}
\usepackage{geometry}
\geometry{margin=2.5cm}

\title{\textbf{Flag Football GPS-Analyse}\\[0.3em]\large Mathematische Ausdrücke}
\author{}
\date{}

\begin{document}
\maketitle

% -------------------------------------------------------
\section*{1. GPS-Konvertierung: Semicircles $\to$ Grad}

Garmin speichert Koordinaten als 32-Bit-Integer (Semicircles):

\[
  \text{lat}_{\deg} = \text{position\_lat}  \times \frac{180}{2^{31}}
  \qquad
  \text{lon}_{\deg} = \text{position\_long} \times \frac{180}{2^{31}}
\]

Ein voller Kreis umfasst $2^{32}$ Semicircles, ein Halbkreis ($180°$) entspricht $2^{31}$.

% -------------------------------------------------------
\section*{2. Geschwindigkeitsglättung (Rolling Mean)}

GPS-Speed-Daten bei 1\,Hz sind verrauscht. Kurze Ausreißer-Peaks werden durch einen
\textbf{zentrierten gleitenden Mittelwert} geglättet: jeder Wert wird durch den
Durchschnitt seiner $w = 3$ Nachbarn (aktueller Punkt + je 1 Sekunde davor/danach) ersetzt.

\[
  \hat{v}(t) = \frac{1}{3} \sum_{k=-1}^{1} v_{\text{raw}}(t + k)
\]

\centering
\textit{Beispiel:} Ein einzelner GPS-Spike auf $8{,}0\ \tfrac{\text{m}}{\text{s}}$
bei sonst $4{,}0\ \tfrac{\text{m}}{\text{s}}$\\ wird auf
$\frac{4{,}0 + 8{,}0 + 4{,}0}{3} = 5{,}3\ \tfrac{\text{m}}{\text{s}}$ reduziert

% -------------------------------------------------------
\section*{4. Sprint-Schwelle}

Die Sprint-Schwelle wird \textbf{relativ zur persönlichen Maximalgeschwindigkeit}
$v_{\max}^{\text{pers}}$ gesetzt (berechnet über alle Sessions). Ein fixer absoluter
Wert (z.\,B.\ $18\ \tfrac{\text{km}}{\text{h}}$) wäre ungeeignet, da Amateure
unterschiedliche Leistungsniveaus haben \cite{sfs_metrics}.

\[
  \theta_{\text{sprint}} = 0{,}60 \cdot v_{\max}^{\text{pers}}
\]

Ein Sample $t$ gilt als Sprint, wenn: \[\hat{v}(t) \;\geq\; \theta_{\text{sprint}}\]


\textit{Beispiel mit $v_{\max}^{\text{pers}} = 6{,}75\ \tfrac{\text{m}}{\text{s}}$
$(= 24{,}3\ \tfrac{\text{km}}{\text{h}})$:}
\[
  \theta_{\text{sprint}} = 0{,}60 \times 6{,}75 = 4{,}05\ \tfrac{\text{m}}{\text{s}}
  \;=\; 14{,}6\ \tfrac{\text{km}}{\text{h}}
\]

Zusätzlich gilt: eine Sprint-Phase muss mindestens $2$ Sekunden andauern
(Mindestdauer-Filter gegen GPS-Rauschen).

% -------------------------------------------------------
\section*{5. Segmentierung (Zustandswechsel)}

Aus der binären Aktivierungsmaske $a(t) \in \{0,1\}$ werden Segmente gebildet:

\[
  \Delta(t) =
  \begin{cases}
    1 & \text{falls } a(t) \neq a(t-1) \\
    0 & \text{sonst}
  \end{cases}
  \qquad
  \text{segment\_id}(t) = \sum_{s=0}^{t} \Delta(s)
\]

% -------------------------------------------------------
\section*{6. Speed-Normalisierung (Heatmap)}

Für jeden Sprint wird die geglättete Geschwindigkeit auf $[0,\,1]$ normiert,
um eine einheitliche Farbkodierung in der GPS-Karte zu ermöglichen
(violett\,=\,langsam, gelb\,=\,schnell):

\[
  \tilde{v}(t) = \frac{\hat{v}(t) - \min(\hat{v})}{\max(\hat{v}) - \min(\hat{v}) + \varepsilon}
  \qquad \varepsilon = 10^{-9}
\]

% -------------------------------------------------------
\section*{7. Bewegungsrichtung (Quiver)}

Die lokale Richtung ergibt sich aus den GPS-Differenzen aufeinanderfolgender Punkte:

\[
  \Delta\lambda(t) = \lambda(t) - \lambda(t-1)
  \qquad
  \Delta\varphi(t) = \varphi(t) - \varphi(t-1)
\]

wobei $\lambda$ die Länge und $\varphi$ die Breite bezeichnet. Der Richtungsvektor
$(\Delta\lambda,\, \Delta\varphi)$ wird alle $n = 3$ Punkte als Pfeil gezeichnet.

\end{document}